% !TeX root = ../main.tex
% Fonte: Lezione 10, p. 1.
% Continua la sezione sulle basi di sottospazi di R^4.

\subsection{Da una base alle equazioni cartesiane}

Una base $u_1,\ldots,u_r$ di $U\subseteq K^m$ fornisce la descrizione \emph{parametrica}
\[
  U=\{a_1u_1+\cdots+a_ru_r:a_1,\ldots,a_r\in K\}.
\]
Gli scalari $a_i$ sono i \emph{parametri}. Le equazioni \emph{cartesiane} impongono condizioni alle sole coordinate: si ottengono eliminando i parametri e verificando che le condizioni siano sia necessarie sia sufficienti.

\begin{esercizio}
Scrivere le equazioni del sottospazio $U\subseteq\R^4$ generato da
\[
  v_1=(1,-1,0,2),\quad v_2=(-1,-3,2,0),\quad v_3=(2,0,-1,3).
\]
\end{esercizio}

\begin{svolgimento}
\textbf{Descrizione parametrica.} Dalla relazione $v_2=3v_1-2v_3$ si usa la base $(v_1,v_3)$. Si ha $x\in U$ se e solo se esistono $a_1,a_2\in\R$ tali che
\[
  x=a_1(1,-1,0,2)+a_2(2,0,-1,3),\qquad
  \begin{cases}
    x_1=a_1+2a_2,\\
    x_2=-a_1,\\
    x_3=-a_2,\\
    x_4=2a_1+3a_2.
  \end{cases}
\]

\textbf{Condizioni necessarie.} La seconda e la terza equazione danno $a_1=-x_2$ e $a_2=-x_3$. Sostituendo nelle altre si ricava il sistema omogeneo
\[
  \begin{cases}
    x_1+x_2+2x_3=0,\\
    2x_2+3x_3+x_4=0.
  \end{cases}
\]

\textbf{Condizioni sufficienti.} Se $x$ soddisfa il sistema, scegliamo $a_1=-x_2$ e $a_2=-x_3$. Allora
\[
  a_1v_1+a_2v_3=(-x_2-2x_3,\ x_2,\ x_3,\ -2x_2-3x_3)
  =(x_1,x_2,x_3,x_4).
\]
Pertanto $x\in U$: il sistema caratterizza esattamente il sottospazio.
\end{svolgimento}

